% Common packages
\usepackage{amsmath}
\usepackage{amssymb}
\usepackage{hyperref}
\usepackage{graphicx}
\usepackage{booktabs}
\usepackage{multirow}
\usepackage{csquotes}
\usepackage{verbatim}
\usepackage{enumitem}
\usepackage{microtype}
\usepackage{tikz}
\usetikzlibrary{arrows.meta, positioning, shapes.geometric, shadows.blur, calc, fit, backgrounds}
\usepackage{fontawesome5}
\usepackage{pgfgantt}
\usepackage{pgfplots}
\usepackage{xstring}
\setlength{\marginparwidth}{2.5cm}
\usepackage[colorinlistoftodos, textwidth=2.3cm]{todonotes}
\usepackage[backend=biber,style=numeric]{biblatex}

% Build system: override output PDF name
\newcommand{\docname}[1]{}

% Common macros
\newcommand{\scare}[1]{\enquote{#1}}
\newcommand{\term}[1]{\textbf{#1}}

\newtheorem{definition}{Definition}

% --- Sticky-note fixme system ---
% Usage: \fixme{text to fix} or \fixme[inline]{longer note}
% When all \fixme commands are removed, the WORKING DRAFT notice no longer appears.
% Comment out the line below to force-hide draft notice even with fixmes present.
\newif\ifdraftmode
\draftmodetrue

\newcommand{\fixme}[2][]{%
  \todo[color=orange!30, size=\small, #1]{#2}%
}
\newcommand{\fixmeinline}[1]{%
  \todo[color=orange!30, inline, size=\small]{#1}%
}

% Shorthand for flagging a missing citation, e.g. \cn after a claim.
% Optional detail: \cn[which reference to find]
\newcommand{\cn}[1][]{%
  \IfStrEq{#1}{}{\fixme{Citation needed.}}{\fixme{Citation needed: #1}}%
}

% Figure convenience: \cfigure[placement][width]{filename}{caption}
% The label is auto-derived from the filename without extension, e.g.
% \cfigure{bootstrap_stability.png}{Caption text} -> \label{fig:bootstrap_stability}
% Images are loaded from the repo's code/figures/ dir (requires the xstring package)
% The optional width defaults to 0.8\textwidth; use [\textwidth] for full width.
% Strip the file extension for use as a label; if there is no dot
% (e.g. a bare stem like "bootstrap_stability") the whole name is used.
% Note: \StrBefore alone would return an empty result when no dot exists.
\newcommand{\stripext}[2]{%
   \IfSubStr{#1}{.}{\StrBefore{#1}{.}[#2]}{\def#2{#1}}%
}

\NewDocumentCommand{\cfigure}{O{htbp} O{0.8\textwidth} m m}{%
   \stripext{#3}{\cfiglabel}%
   \begin{figure}[#1]
      \centering
      \includegraphics[width=#2]{../../code/figures/#3}
      \caption{#4}
      \label{fig:\cfiglabel}
   \end{figure}%
}

% Figure reference convenience: \figref{filename} -> "Figure N"
% Uses the same filename-derived label as \cfigure, e.g.
% \cfigure{bootstrap_stability.png}{...} then \figref{bootstrap_stability.png}
\newcommand{\figref}[1]{%
   \stripext{#1}{\cfiglabel}%
   Figure~\ref{fig:\cfiglabel}%
}

% Table convenience: \ctable[placement]{filename}{caption}
% The label is auto-derived from the filename without extension, e.g.
% \ctable{fine-tuning/main-table-Qwen3.5-9B.tex}{...} -> \label{tab:main-table-Qwen3.5-9B}
% The file is expected to hold a bare tabular (e.g. pandas to_latex output)
% and is loaded from the repo's code/figures/ dir (requires the xstring package)
\newcommand{\ctable}[3][htbp]{%
   \stripext{#2}{\ctablabel}%
   \begin{table}[#1]%
      \centering\small%
      \input{../../code/figures/#2}%
      \caption{#3}%
      \label{tab:\ctablabel}%
   \end{table}%
}

% Table reference convenience: \tabref{filename} -> "Table N"
% Uses the same filename-derived label as \ctable, e.g.
% \ctable{fine-tuning/main-table-Qwen3.5-9B.tex}{...} then \tabref{fine-tuning/main-table-Qwen3.5-9B.tex}
\newcommand{\tabref}[1]{%
   \stripext{#1}{\ctablabel}%
   Table~\ref{tab:\ctablabel}%
}

% --- Research question (RQ) system ---
% Define the research questions once (e.g. in the introduction) and then
% cross-reference any of them throughout the document with a clickable link.
% In an electronic PDF the reference shows a tooltip with the full question
% text when hovered; in print only the "RQn" link label is seen. The tooltip
% is a plain PDF Widget annotation (no extra packages required) - it works in
% desktop readers (Okular, Evince, Acrobat); some viewers ignore it.
%
% Usage (the label string needs no "rq:" prefix - it is added automatically):
%   \newrq{label}{Question text ...}  % define & print an RQ (auto-numbered)
%   \rqref{label}                     % "RQn" clickable link (hover for full text)
%   \rqtext{label}                    % the full question text inline in prose
%
% Example:
%   \newrq{align}{Can we ...?}
%   ... see \rqref{align} in Chapter 3.
\newcounter{rqcount}
\newcommand{\newrq}[2]{%
   \stepcounter{rqcount}%
   \expandafter\xdef\csname rqnum@#1\endcsname{RQ\number\csname c@rqcount\endcsname}%
   \expandafter\gdef\csname rqtext@#1\endcsname{#2}%
   \hypertarget{rq:#1}{}%
   \textbf{\csname rqnum@#1\endcsname.}\ #2%
}
% \pdfhover{visible text}{hover text} - a hover tooltip annotation overlaying
% the visible text (PDF Widget button with zero border). Used by the RQ and
% glossary reference macros below. The hover text is fully expanded into
% \tttmp first so \pdfescapestring sees the literal string (safe for
% parentheses/quotes). The dict is fully expanded at call time (hence
% \expanded), so several tooltips on one page each keep their own text.
\newcommand{\pdfhover}[2]{%
   \edef\tttmp{#2}%
   \leavevmode
   \setbox0=\hbox{#1}%
   \expanded{\noexpand\pdfannot width \the\wd0 height \the\ht0 depth \the\dp0 %
      {/Subtype /Widget /FT /Btn /Ff 65536 /F 768 /H /N /BS << /W 0 >>
       /TU (\pdfescapestring{\tttmp}) /Contents (\pdfescapestring{\tttmp})}}%
   #1%
}
% \rqref and \termref use \texorpdfstring so bookmarks/PDF outlines get plain
% text, and \protect so the links survive writes (e.g. \section headings and
% \caption arguments going into the TOC/LOF): the write keeps the tokens and
% the link is typeset when the list is rendered.
\newcommand{\rqref}[1]{%
   \texorpdfstring{\protect\hyperlink{rq:#1}{\protect\pdfhover{\csname rqnum@#1\endcsname}{\csname rqtext@#1\endcsname}}}{\csname rqnum@#1\endcsname}%
}
\newcommand{\rqtext}[1]{\csname rqtext@#1\endcsname}

% --- Glossary of terms system ---
% Define terms once (e.g. in a "Glossary of terms" section in the
% introduction) and then reference them throughout the document with a
% clickable link. In an electronic PDF the reference shows the definition as
% a tooltip when hovered (same mechanism as the RQ references); in print only
% the linked term text is seen.
%
% Usage (labels need no prefix - internal namespacing is automatic):
%   \newterm{label}{Term text}{Definition text}  % register a term
%   \termref{label}              % clickable term link (hover tooltip)
%   \termtext{label}             % the definition text inline in prose
%   \printglossary               % prints the entries only: wrap it in your
%                                % own \section / \chapter heading, e.g. an
%                                % appendix chapter.
%
% Example:
%   \newterm{llm}{Large language model}{A neural network trained on ...}
%   ... see \termref{llm} ...
%
% Terms must be registered (\newterm) before their first reference, the same
% restriction as for research questions. In the report the terms live in a
% single file (appendices/glossary.tex) which is \input twice: once in the
% preamble (registers the terms so \termref works anywhere) and once as the
% appendix (prints them; set \glossaryprinttrue first). \newterm is
% idempotent, so the second \input does not register the terms twice.
\newif\ifglossaryprint
\newcommand{\glistall}{}
\newcommand{\newterm}[3]{%
   \ifcsname gldef-#1\endcsname
   \else
      \expandafter\gdef\csname gltxt-#1\endcsname{#2}%
      \expandafter\gdef\csname gldef-#1\endcsname{#3}%
      \xdef\glistall{\ifx\glistall\empty\else\glistall,\fi#1}%
   \fi
}
\newcommand{\termref}[1]{%
   \texorpdfstring{\protect\hyperlink{gl:#1}{\protect\pdfhover{\csname gltxt-#1\endcsname}{\csname gldef-#1\endcsname}}}{\csname gltxt-#1\endcsname}%
}
\newcommand{\termtext}[1]{\csname gldef-#1\endcsname}
\makeatletter
\newcommand{\printglossary}{%
   \expandafter\gl@printlist\expandafter{\glistall}%
}
\newcommand{\gl@printlist}[1]{%
   \@for\gl@lab:=#1\do{%
      \par\noindent
      \hypertarget{gl:\gl@lab}{}\textbf{\csname gltxt-\gl@lab\endcsname:}\ %
      \csname gldef-\gl@lab\endcsname\par\vspace{6pt}%
   }%
}
\makeatother
